% !PRINTED 356
\weberpage{356}
%<<CONT
solcher Gleichungen erhalten wir aber noch zwei, wenn wir (1) mit $x^{-m}$ und mit $x^{-m-n}$ multipliciren. Dies giebt
\[
(3)\quad r^{n+m} = \frac{-g\sin m\vartheta}{\sin n\vartheta},
\]
\[
(4)\quad r^m = \frac{-g\sin(n+m)\vartheta}{e\sin n\vartheta},
\]
und von diesen drei Gleichungen folgt jede aus den beiden anderen. Wenn man $r$ aus zwei von ihnen eliminirt, so ergiebt sich, wenn, wie oben
\[
\lambda = \frac{g^n}{e^{n+m}}
\]
gesetzt wird,
\[
(5)\quad \lambda = (-1)^n \frac{\sin n\vartheta^n \sin m\vartheta^m}{\sin(n+m)\vartheta^{n+m}}.
\]

Man kann, da man von den beiden conjugirten Wurzeln nur die eine zu berechnen braucht, $\vartheta$ auf den ersten Quadranten beschränken, muss aber dann unter Umständen auch $r$ negativ annehmen.

Wie man nun aus der Formel (5) mittelst der trigonometrischen Tafeln den Winkel $\vartheta$ und dann aus einer der drei Formeln (2), (3), (4) den zugehörigen Werth von $r$ berechnet, das wollen wir nun an dem vorhin betrachteten Beispiel
\[
x^3 - 2x - 2 = 0
\]
noch nachweisen.

Hier wird die Gleichung (5)
\[
(6)\quad \frac{1}{2} = \frac{\sin 2\vartheta^2 \sin\vartheta}{\sin 3\vartheta^3},
\]
oder
\[
(7)\quad 3\log\sin 3\vartheta - 2\log\sin 2\vartheta - \log\sin\vartheta = 0{,}3010300,
\]
woraus zunächst ersichtlich ist, dass der Winkel $\vartheta$ in dem Intervall von 0 bis $60^0$ liegt. Man findet zunächst nach wenigen Versuchen, dass $\vartheta$ zwischen $33^0$ und $34^0$ liegt, und wenn man dann für diese beiden Werthe die Differenz
\[
(8)\quad 3\log\sin 3\vartheta - 2\log\sin 2\vartheta - \log\sin\vartheta - 0{,}3010300
\]
berechnet, so erhält man zunächst auf drei Decimalen
\[
0{,}026, \quad -0{,}013,
\]
und hieraus ergiebt sich durch Interpolation der genauere Werth
\[
\vartheta = 33^0\,40'.
\]
